\textbf{Proof.}
Let \(k=2\lfloor d/2\rfloor\).  Because \(2\leq d\leq n^2\),
\[
 0<\alpha^2=\frac{\sqrt{k}}{16n}\leq\frac1{16},
\]
so \(0<\alpha\leq1/4\), every displayed coordinate of \(\mathbf P_\sigma\) is nonnegative, each perturbed pair sums to \(2/k\), and \(\mathbf P_\sigma\in\Delta_d\).  Directly,
\[
 \lVert\mathbf P_\sigma-\mathbf Q\rVert_1
 =\sum_{j=1}^{k/2}\left(\frac\alpha k+\frac\alpha k\right)
 =\alpha.
\]

Let \(M=2^{-k/2}\sum_\sigma\mathbf P_\sigma^n\).  For two independent uniform sign vectors \(\sigma,\tau\), the standard likelihood-ratio second-moment identity, evaluated coordinate by coordinate under \(\mathbf Q\), gives
\[
 \begin{aligned}
 1+\chi^2(M,\mathbf Q^n)
 &=E_{\sigma,\tau}
 \left(\sum_{x=1}^d\frac{P_{\sigma,x}P_{\tau,x}}{Q_x}\right)^n\\
 &=E_{\sigma,\tau}\left[
 1+\frac{2\alpha^2}{k}
      \sum_{j=1}^{k/2}\sigma_j\tau_j
 \right]^n\\
 &\leq E_{\sigma,\tau}\exp\!\left{
  \frac{2n\alpha^2}{k}
      \sum_{j=1}^{k/2}\sigma_j\tau_j
 \right}\\
 &=\left\{\cosh\!\left(\frac{2n\alpha^2}{k}\right)\right\}^{k/2}
 \leq\exp\!\left(\frac{n^2\alpha^4}{k}\right)
 =e^{1/256}.
 \end{aligned}
\]
Coordinates with \(Q_x=0\) also have \(P_{\sigma,x}=0\) and may be omitted from the likelihood ratio.  The two inequalities use \(1+t\leq e^t\) and \(\cosh t\leq e^{t^2/2}\).  This proves the stated chi-square bound.

For completeness, convert the fuzzy hypotheses to squared-error risk with their constants.  Since
\[
 \lVert M-\mathbf Q^n\rVert_{\mathrm{TV}}
 \leq\frac12\sqrt{\chi^2(M,\mathbf Q^n)}
 \leq\frac12\sqrt{e^{1/256}-1}<\frac12,
\]
the test \(\phi=\mathbf1\{\widehat L\geq\alpha/2\}\) satisfies
\[
 \begin{aligned}
 \sup_{(\mathbf P,\mathbf Q)\in\Delta_d^2}
 E[(\widehat L-\lVert\mathbf P-\mathbf Q\rVert_1)^2]
 &\geq\frac{\alpha^2}{8}
 \{\mathbf Q^n(\phi=1)+M(\phi=0)\}\\
 &\geq\frac{\alpha^2}{8}
 \{1-\lVert M-\mathbf Q^n\rVert_{\mathrm{TV}}\}
 \geq\frac{\alpha^2}{16}.
 \end{aligned}
\]
The independent \(\mathbf Q\)-sample has the same distribution under the null and every mixture component and therefore cannot reduce the testing affinity.  Taking the infimum over estimators yields
\[
 \mathfrak R^{L_1}_{n,d}\geq\frac{\alpha^2}{16}
 =\frac{\sqrt{k}}{256n}.
\]
Finally, \(\mathrm{lem:l1-lower-transfer}\) gives
\[
 \mathfrak R_{n,d,\epsilon}
 \geq\frac1{16}\mathfrak R^{L_1}_{n,d}
 \geq\frac{\sqrt{k}}{4096n}
 \geq\frac{\sqrt d}{4096\sqrt2\,n},
\]
where \(k=2\lfloor d/2\rfloor\geq d/2\).  The embedded causal laws have propensity \(1/2\), which obeys fixed overlap for every \(0<\epsilon<1/2\), including all boundary and null cells specified in \(\mathrm{def:l1-embedding}\).